\scanpage{10}
\handwritten{24.4,67}\\

Dear Murre,\\

I am sending you enc\unsure{\deleted{o}l|\deleted{l}o}osed the sketch which I promised on ramificatio\unsure{n}
% The terminal character has a surviving left stem/serif; its right side is cut off by the physical right edge.
data, as well as your manuscript together with my annotations. As the paragraphs 1 to 5 will have to be changed entirely, there is maybe not much point for you to look up to carefully my annotations on these. \deleted{As} I would only suggest that you put more fo\unsure{\deleted{m}r|\deleted{r}m|r}mulas, arrows etc on \deleted{separ\unsure{a|t}te} separate lines, for easier reading; that you normalize the sta\unsure{\deleted{n}d|\deleted{d}n|d}\unsure{\deleted{d}f|\deleted{f}d|f}ard references to SGA, EGA (for instance, no need to indicate SGA when reference is being made to the seminar 1960-61 in which your talk is going to be included itself, reference as (XI 6.1) is then sufficient); that you be cautious when using a locution like ``if primitive $n$.th roots exist over $S$'', which does not mean much by itself, unless you make clear what you really mean. I would propose the following tentative plan for your talk:
% In the three crowded words above, alternative overprint orders are not treated as established correction chronology.

\begin{enumerate}[label={}]
\item[1.] Kummer coverings with respect to a given set of sections.
\item[2.] Ramification data.
\reversemarginpar
\margin{\handwritten{\unsure{\figure{\begin{tikzpicture}[x=.014cm,y=.014cm]
\draw[gray,line width=.25pt] (26,64) .. controls (29,51) and (34,35) .. (30,22) .. controls (28,10) and (14,4) .. (0,0);
\end{tikzpicture}}}}}
\normalmarginpar
% Faint hooked trace in the left margin, between the levels of items 2 and 3; intentional mark versus incidental trace remains uncertain.
\item[3.] Properties of the catégory of coverings of type $\underline{R}$. The group ${}_{1}^{\underline{R}}(S,\ )$
% No base glyph is visible before the stacked R and 1. The second argument is blank. The closing parenthesis meets the paper edge; no following punctuation is visible.
\item[4.] Functorial behaviours.
\item[5.] Galois coverings of type $\underline{R}$.
\item[6.] Abhyankars theorem (for divisors with normal crossings).
\end{enumerate}
etc

The sections 6 to 10 essentially the same as in your notes, just taking into account my observations, and changing inner references according to the changes made in 1. to 5. A paragraph 11 could be an appendix on the tame variant of Lefschetz's theorem. - Par.1 is essentially your previous par.1 ; you should check exactly what is needed in it for \unsure{\deleted{f}w|\deleted{w}f}hat
% archive sheet number: 9
